% GENERATED FILE. Edit canonical lessons, then run npm run generate:books.
% canonical-source-hash: fnv1a64:4dcda2104198ac64
% canonical-lessons: TA-C366-turai, TA-C366-ulaviyal, TA-C366-puviyiyal, TA-C366-cuttiram, TA-C366-varaiyaru

\chapter{Fields of Knowledge}
\label{ch:ta-fields-of-knowledge}
\hlchaptermodality{366}

\begin{chapteropening}
\textbf{By the end of this chapter:} \emph{I can name some fields of study and define a word.}

Complete the last lesson of chapter 366: I can name some fields of study and define a word.
\end{chapteropening}

\section[\texorpdfstring{\ta{துறை} (tuṟai)}{துறை (tuṟai)}]{\ta{துறை} (tuṟai) — a field; a department}

\label{lesson:TA-C366-turai}



\begin{quote}
Before the new one: say the Tamil for a summary, then the Tamil for a quotation.
\end{quote}

\subsection*{You'll want to know: \ta{துறை}}
\textbf{\ta{துறை}} — \emph{tuṟai} — \textquotedblleft{}a field; a department\textquotedblright{}.

\begin{quote}
\ta{நீங்கள்} \ta{எந்தத்} \ta{துறையில்} \ta{வேலை} \ta{செய்கிறீர்கள்}?
\end{quote}

\begin{quote}
\emph{nīṅkaḷ entat tuṟaiyil vēlai ceykiṟīrkaḷ?}
\end{quote}

\begin{quote}
Which field do you work in?
\end{quote}

The same word names a landing place on a river; \ta{துறைமுகம்}, a harbour, is built on it.

\begin{cousinweb}
\textbf{\ta{துறை}} (\emph{tuṟai}) is native Dravidian: DEDR 3370 lists Tamil \emph{tuṟai}, \textquotedblleft{}place, way, branch, section, category; landing place, harbour\textquotedblright{}.
\end{cousinweb}

\begin{grammarlens}[title={What you've built}]
One more word for giving reasons and opinions.
\end{grammarlens}

\subsection*{Your turn}
\begin{itemize}
\raggedright
  \item \emph{Say it:} \emph{tuṟai}
  \item \emph{Say it:} \emph{tuṟai}, once more
  \item \emph{Say it:} \emph{nīṅkaḷ entat tuṟaiyil vēlai ceykiṟīrkaḷ?} — Which field do you work in?
  \item \emph{Recall:} say the Tamil for a summary, then the Tamil for a quotation, then say \emph{tuṟai} again
\end{itemize}

\begin{tcolorbox}[breakable,colback=teal!4,colframe=teal!35!black,title={Before you move on}]
What does \ta{துறை} mean? (\textquotedblleft{}A field; a department\textquotedblright{}.) Say it once more.
\end{tcolorbox}
\section[\texorpdfstring{\ta{உளவியல்} (uḷaviyal)}{உளவியல் (uḷaviyal)}]{\ta{உளவியல்} (uḷaviyal) — psychology}

\label{lesson:TA-C366-ulaviyal}



\begin{quote}
Before the new one: say the Tamil for a quotation, then the Tamil for a field.
\end{quote}

\subsection*{You'll want to know: \ta{உளவியல்}}
\textbf{\ta{உளவியல்}} — \emph{uḷaviyal} — \textquotedblleft{}psychology\textquotedblright{}.

\begin{quote}
\ta{அவள்} \ta{பல்கலைக்கழகத்தில்} \ta{உளவியல்} \ta{படிக்கிறாள்}.
\end{quote}

\begin{quote}
\emph{avaḷ palkalaikkaḻakattil uḷaviyal paṭikkiṟāḷ.}
\end{quote}

\begin{quote}
She studies psychology at university.
\end{quote}

\ta{உளம்}, the mind + -\ta{இயல்}, a field of study, as in \ta{அறிவியல்}, science.

\begin{cousinweb}
\textbf{\ta{உளவியல்}} (\emph{uḷaviyal}) is a modern native coinage: \emph{uḷam}, \textquotedblleft{}mind, heart\textquotedblright{}, in the family of \emph{uḷ}, \textquotedblleft{}inside\textquotedblright{} (DEDR 698), and \emph{iyal}, \textquotedblleft{}field of study\textquotedblright{} (DEDR 471).
\end{cousinweb}

\begin{grammarlens}[title={What you've built}]
One more word for giving reasons and opinions.
\end{grammarlens}

\subsection*{Your turn}
\begin{itemize}
\raggedright
  \item \emph{Say it:} \emph{uḷaviyal}
  \item \emph{Say it:} \emph{uḷaviyal}, once more
  \item \emph{Say it:} \emph{avaḷ palkalaikkaḻakattil uḷaviyal paṭikkiṟāḷ.} — She studies psychology at university.
  \item \emph{Recall:} say the Tamil for a quotation, then the Tamil for a field, then say \emph{uḷaviyal} again
\end{itemize}

\begin{tcolorbox}[breakable,colback=teal!4,colframe=teal!35!black,title={Before you move on}]
What does \ta{உளவியல்} mean? (\textquotedblleft{}Psychology\textquotedblright{}.) Say it once more.
\end{tcolorbox}
\section[\texorpdfstring{\ta{புவியியல்} (puviyiyal)}{புவியியல் (puviyiyal)}]{\ta{புவியியல்} (puviyiyal) — geography}

\label{lesson:TA-C366-puviyiyal}



\begin{quote}
Before the new one: say the Tamil for a field, then the Tamil for psychology.
\end{quote}

\subsection*{You'll want to know: \ta{புவியியல்}}
\textbf{\ta{புவியியல்}} — \emph{puviyiyal} — \textquotedblleft{}geography\textquotedblright{}.

\begin{quote}
\ta{புவியியல்} \ta{வகுப்பில்} \ta{நதிகளைப்} \ta{பற்றிப்} \ta{படித்தோம்}.
\end{quote}

\begin{quote}
\emph{puviyiyal vakuppil natikaḷaip paṟṟip paṭittōm.}
\end{quote}

\begin{quote}
In geography class we studied rivers.
\end{quote}

\ta{புவி}, the earth + -\ta{இயல்}, the study of.

\begin{cousinweb}
\textbf{\ta{புவியியல்}} (\emph{puviyiyal}) joins \emph{puvi}, \textquotedblleft{}the earth\textquotedblright{}, borrowed from Sanskrit (compare \emph{bhū}, \textquotedblleft{}the earth\textquotedblright{}, Monier-Williams), and native \emph{iyal}, \textquotedblleft{}field of study\textquotedblright{} (DEDR 471).
\end{cousinweb}

\begin{grammarlens}[title={What you've built}]
One more word for giving reasons and opinions.
\end{grammarlens}

\subsection*{Your turn}
\begin{itemize}
\raggedright
  \item \emph{Say it:} \emph{puviyiyal}
  \item \emph{Say it:} \emph{puviyiyal}, once more
  \item \emph{Say it:} \emph{puviyiyal vakuppil natikaḷaip paṟṟip paṭittōm.} — In geography class we studied rivers.
  \item \emph{Recall:} say the Tamil for a field, then the Tamil for psychology, then say \emph{puviyiyal} again
\end{itemize}

\begin{tcolorbox}[breakable,colback=teal!4,colframe=teal!35!black,title={Before you move on}]
What does \ta{புவியியல்} mean? (\textquotedblleft{}Geography\textquotedblright{}.) Say it once more.
\end{tcolorbox}
\section[\texorpdfstring{\ta{சூத்திரம்} (cūttiram)}{சூத்திரம் (cūttiram)}]{\ta{சூத்திரம்} (cūttiram) — a formula}

\label{lesson:TA-C366-cuttiram}



\begin{quote}
Before the new one: say the Tamil for psychology, then the Tamil for geography.
\end{quote}

\subsection*{You'll want to know: \ta{சூத்திரம்}}
\textbf{\ta{சூத்திரம்}} — \emph{cūttiram} — \textquotedblleft{}a formula\textquotedblright{}.

\begin{quote}
\ta{ஆசிரியர்} \ta{கரும்பலகையில்} \ta{ஒரு} \ta{சூத்திரத்தை} \ta{எழுதினார்}.
\end{quote}

\begin{quote}
\emph{āciriyar karumpalakaiyil oru cūttirattai eḻutiṉār.}
\end{quote}

\begin{quote}
The teacher wrote a formula on the blackboard.
\end{quote}

Sanskrit sūtra is a thread: a short rule that strings ideas together. \ta{வெற்றியின்} \ta{சூத்திரம்}, the formula for success.

\begin{cousinweb}
\textbf{\ta{சூத்திரம்}} (\emph{cūttiram}) is borrowed from Sanskrit \emph{sūtra}, \textquotedblleft{}thread, cord\textquotedblright{} (CDIAL 13561), which also named a short rule in grammar and science. Tamil writes the s as \ta{ச}.
\end{cousinweb}

\begin{grammarlens}[title={What you've built}]
One more word for giving reasons and opinions.
\end{grammarlens}

\subsection*{Your turn}
\begin{itemize}
\raggedright
  \item \emph{Say it:} \emph{cūttiram}
  \item \emph{Say it:} \emph{cūttiram}, once more
  \item \emph{Say it:} \emph{āciriyar karumpalakaiyil oru cūttirattai eḻutiṉār.} — The teacher wrote a formula on the blackboard.
  \item \emph{Recall:} say the Tamil for psychology, then the Tamil for geography, then say \emph{cūttiram} again
\end{itemize}

\begin{tcolorbox}[breakable,colback=teal!4,colframe=teal!35!black,title={Before you move on}]
What does \ta{சூத்திரம்} mean? (\textquotedblleft{}A formula\textquotedblright{}.) Say it once more.
\end{tcolorbox}
\section[\texorpdfstring{\ta{வரையறு} (varaiyaṟu)}{வரையறு (varaiyaṟu)}]{\ta{வரையறு} (varaiyaṟu) — to define}

\label{lesson:TA-C366-varaiyaru}



\begin{quote}
Before the new one: say the Tamil for geography, then the Tamil for a formula.
\end{quote}

\subsection*{You'll want to know: \ta{வரையறு}}
\textbf{\ta{வரையறு}} — \emph{varaiyaṟu} — \textquotedblleft{}to define\textquotedblright{}.

\begin{quote}
\ta{இந்தச்} \ta{சொல்லின்} \ta{பொருளை} \ta{அவர்} \ta{தெளிவாக} \ta{வரையறுத்தார்}.
\end{quote}

\begin{quote}
\emph{intac colliṉ poruḷai avar teḷivāka varaiyaṟuttār.}
\end{quote}

\begin{quote}
He clearly defined the meaning of this word.
\end{quote}

Past (I …): \ta{வரையறுத்தேன்}, \emph{varaiyaṟuttēṉ}.

What is defined takes -\ta{ஐ}. The noun is \ta{வரையறை}, a definition.

\begin{cousinweb}
\textbf{\ta{வரையறு}} (\emph{varaiyaṟu}) is native: \emph{varai}, \textquotedblleft{}boundary, limit\textquotedblright{} (DEDR 5261), and \emph{aṟu}, \textquotedblleft{}to be severed, cut\textquotedblright{} (DEDR 315): to cut out the boundary of a meaning.
\end{cousinweb}

\begin{grammarlens}[title={What you've built}]
One more word for giving reasons and opinions.
\end{grammarlens}

\subsection*{Your turn}
\begin{itemize}
\raggedright
  \item \emph{Say it:} \emph{varaiyaṟu}
  \item \emph{Say it:} \emph{varaiyaṟu}, once more
  \item \emph{Say it:} \emph{intac colliṉ poruḷai avar teḷivāka varaiyaṟuttār.} — He clearly defined the meaning of this word.
  \item \emph{Recall:} say the Tamil for geography, then the Tamil for a formula, then say \emph{varaiyaṟu} again
\end{itemize}

\begin{tcolorbox}[breakable,colback=teal!4,colframe=teal!35!black,title={Before you move on}]
What does \ta{வரையறு} mean? (\textquotedblleft{}To define\textquotedblright{}.) Say it once more.
\end{tcolorbox}
